\documentclass[10pt]{article}

\usepackage[utf8]{inputenc}

\usepackage[T1]{fontenc}

\usepackage{amsmath}

\usepackage{amsfonts}

\usepackage{amssymb}

\usepackage[version=4]{mhchem}

\usepackage{stmaryrd}



\begin{document}

является аналитич. функцией в области, не содержащей нулей статистич. суммы, то есть при $|z| \neq 1$. Теоремы Ли Янга справедливы и в случае двух-, трехкомпонентных векторнозначных спинов (см. [5]).



Jum.: [1] Yang C. N., Le e T. D., «Phys. Rev.», 1952, v. 87, p. 404; там же, р. 410; [2] Рюэл ь Д., Статистическая механика. Строгие результаты, пер. с англ., М., 1971; [3] Ruelle D., «Comm. Math. Phys.», 1970, v. 18, p. 127; [4] Майер Дж., ГеппертМ ай е р М., Статистическая механика, пер. с англ., М., 1980; [5] Гл и м м Дж., Джа ф фе А., Математические методы квантовой физики, пер. с англ., М., 1984.



П. В. Мальиев.



ЛИДСКОГО ТЕОРЕМА - см. След оператора.



ЛИНДБЛАДА ФОРМУЛА - см. Динамическая полугруппа. ЛИНДЕЛЁФА ПРИНЦИП - основной качественныЙ вариационный принцип в теории конформного отображения, найденный Э. Линделёфом (E. Lindelöf, 1915). Пусть односвязные области $D$ и $\widetilde{D}$ на плоскости комплексного переменного $z$ таковы, что их границы $\Gamma$ и $\widetilde{\Gamma}$ соответственно состоят из конечного числа жордановых дуг, причем $\tilde{D}$ содержится в $D$, и пусть точка $z_{0} \in \widetilde{D} \subset D$. Пусть, кроме того, $w=f(z)$ и $w=\tilde{f}(z)-$ функщии, реализующие конформное отображение соответственно $D$ и $\widetilde{D}$ на единичный круг $\Delta=\{w:|w|<1\}$, причем $f\left(z_{0}\right)=0, \tilde{f}\left(z_{0}\right)=0$. П р и нци п Линделё $ф$ а состоит в том, что при этих условиях: 1) прообраз $\widetilde{D}_{\rho}$ области $|w|<\rho, 0<\rho<1$, при отображении $w=\tilde{f}(z)$ лежит внутри прообраза $D_{\rho}$ этой же области $|w|<\rho$ при отображении $w=f(z)$, причем соприкосновение их границ $\widetilde{\Gamma}_{\rho}$ и $\Gamma_{\rho}$ возможно лишь, если $\left.\widetilde{D}=D ; 2\right)\left|\widetilde{f}^{\prime}\left(z_{0}\right)\right| \geqslant\left|f^{\prime}\left(z_{0}\right)\right|$, причем знак равенства возможен лишь, если $\widetilde{D}=D ; 3)$ если существует общая точка $z_{1}$ контуров $\widetilde{\Gamma}$ и $\Gamma$, то



\[

\left|\tilde{f^{\prime}}\left(z_{1}\right)\right| \leqslant\left|f^{\prime}\left(z_{1}\right)\right| \text {, }

\]



причем знак равенства возможен лишь, если $\widetilde{D}=D$. Иными словами, при вдавливании внутрь границы $\Gamma$ области $D$ : 1) все лини и уровня $\Gamma_{\rho}$, то есть прообразы окружностей $|w|=\rho$, сжимаются; 2) растяжение в точке $z_{0}$ увеличивается; 3) растяжение в неподвижных точках $z_{1}$ границы уменышается.



Из приведенной формулировки Л. п. вытекает также, что длина образа дуги $\gamma$ границы $\Gamma$, подвергшейся вдавливанию до дуги $\tilde{\gamma}$, всегда не превосходит длины образа $\tilde{\gamma}$ (длина $f(\gamma) \leqslant$ длины $\tilde{f}(\tilde{\gamma})$ ), причем равенство имеет место только в случае $\tilde{\Gamma}=\Gamma$. Это следствие Л. п. известно также как принцип Монтеля.



Л. п. позволяет получить многие количественные оценки изменения конформного отображения при вариации области. Е.Д. Соломенцев. ЛИНДШТЕЙНА МЕТОД - см. Усреднения метод. ЛИНЕАРИЗОВАННОЕ УРАВНЕНИЕ БОЛЬЦМАНА линейное интегро-дифференциальное уравнение, приближенно описывающее временную эволюцию одночастичной функции распределения разреженного газа вблизи равновесия. Л.у. Б.



\[

\frac{\partial \varphi}{\partial t}+\left(\vec{p} \circ \nabla_{q} \varphi\right)=J_{0}(\varphi) \frac{1}{\varepsilon}

\]



получается из Больцмана уравнения



путем замены



\[

\frac{\partial f}{\partial t}+\left(\vec{p} \circ \nabla_{\vec{q}} f\right)=\frac{1}{\varepsilon} \dot{J}(f, f)

\]



\[

f=\pi^{-3 / 2} \exp \left\{-\vec{p}^{2}\right\}+\mu \exp \left\{-\frac{\vec{p}^{2}}{2}\right\} \varphi

\]



и приравнивания членов, содержащих первую степень параметра $\mu$. Оператор $\dot{J}_{0}$ называется л и не а р и з о в а н н ы м оператором столкновений.



Для Л.у. Б. существует строгое доказательство теоремы сушествования и единственности решения задачи Коши (см. [2]). Лит.: [1] К а рл е м а н Т., Математические задачи кинетической теории газов, пер. с франц., М., 1960; [2] А р с е н ь е в А. А., «Журнал вычислит. математики и математич. физики», 1965, т. 5, № 5, с. 864-82. $\quad$ M. A. Сташенко.



ЛИНЕЙНАЯ ГРУППА - см. Ли групnа.



ЛИНЕЙНАЯ КОМБИНАЦИЯ в е то о ов - см. Векторное пространство.



ЛИНЕЙНАЯ КРАЕВАЯ ЗАДАЧА - см. КраеваЯ Задача.



ЛИНЕЙНАЯ ОБОЛОЧКА Подмножества линейного пространства - совокупность всех линейных комбинаций элементов этого подмножества.



ЛИНЕЙНАЯ СВЯЗНОСТЬ - см. Связность.



ЛИНЕЙНАЯ ЭКСПОНЕНТА - ФУНКЦИЯ вИДа



\[

e^{i z}(\xi)=e^{i z \xi}

\]



где $\xi$ - вещественный или комплексный аргумент. См. Лаnласа преобразование обобщенной функции.



ЛИНЕЙНО НЕЗАВИСИМЫЕ ВЕКТОРЫ - СМ. Векторное пространство.



ЛИНЕЙНО СВЯЗНОЕ ПРОСТРАНСТВО - сМ. СвязНое пространство.



ЛИНЕЙНОТО СОПРЯХЕНИЯ ЗАДАЧА - сМ. КраевЫе задачи теории аналитических функщий.



ЛИНЕЙНОЕ ДИФФЕРЕНЦИАЛЬНОЕ ВЫРАХЕНИЕ см. Дифференциальный оператор.



ЛИНЕЙНОЕ ИНТЕТРАЛЬНОЕ УРАВНЕНИЕ - СМ. Интегральное уравнение.



ЛИНЕЙНОЕ ПРЕДСТАВЛЕНИЕ - см. Представление группы.



ЛИНЕЙНОЕ ПРОСТРАНСТВО - см. Векторное пространство.



ЛИНЕЙНЫЕ КОЛЕБАНИЯ - СМ. МалЫе КолебанИЯ. ЛИНЕЙНЫЙ ГАРМОНИЧЕСКИЙ ОСЦИЛЛЯТОР см. Осииляттор.



ЛИНЕЙНЫЙ ОПЕРАТОР $A$ в векторном пространстве $L-$ отображение, сопоставляющее каждому вектору $e$ некоторого множества $D$ (содержащегося в $L$ и называемого о 6 ла с т ью о п р е де л ен и я Л. о.) другой вектор, обозначаемый $A e$ (и называемый знач е н и м Л.о. на векторе $e$ ). Выполнены следующие условия: 1) $D$-линейное множество, то есть для любых его элементов $e_{1}$ и $e_{2}$ и любых комплексных чисел $\lambda_{1}$ и $\lambda_{2}$ вектор $\lambda_{1} e_{1}+\lambda_{2} e_{2}$ также принадлежит $D$; 2) Л. о. переводит линейную комбинацию векторов в ту же линейную комбинацию соответствующих значений: $A\left(\lambda_{1} e_{1}+\lambda_{2} e_{2}\right)=\lambda_{1} A e_{1}+\lambda_{2} A e_{2}$.



П р и м е р ы Л. о.: матрица



\[

A=\left\|a_{i j}\right\|

\]



действующая в $n$-мерном евклидовом пространстве по правилу



\[

(A e)_{i}=\sum_{j=1}^{n} a_{i j} e_{j},

\]



где векторы $e_{j}$ - столбцы комплексных чисел;



дифференциальный оператор



\[

A=a_{0}+a_{1} d / d x+\ldots+a_{n} d^{n} / d x^{n},

\]



определенный равенством



$(A f)(x)=a_{0} f(x)+a_{1} d f(x) / d x+\ldots+a_{n} d^{n} f(x) / d x^{n} ;$



интегральный оператор $A$, определенный соотношением



\[

(A f)(x)=\int d y A(x, y) f(y) .

\]



Конечномерные пространства. В конечномерном пространстве Л.о. можно определить на всех векторах и задать нек-рой матрицей $\left\|a_{i j}\right\|$. Если $e_{1}, e_{2}, \ldots, e_{n}$ - ортонормированный базис, то $a_{i j}=\left(e_{i}, A e_{j}\right)$, где $(e, f)-$ скалярное произведение векторов $e$ и $f$. Если для нек-рого вектора $e$ и





\end{document}